\section{Heat kernel expansion on manifolds without boundary}
\label{sec4}
\subsection{General formulae}\label{s4gen}
For physicists, the most familiar way to calculate the heat
kernel coefficients is the DeWitt iterative procedure which
will be briefly described in sec.\ \ref{s4iter}. Here
we take a different route following the method of Gilkey.
Advantages of this method are most clearly seen on manifolds
with boundaries. However, even if boundaries are absent
we use the same method for the following reasons:
(i) this will ensure a smooth transition to more
complicated material of the next section; (ii) we believe
that calculations of the coefficients $a_k(f,D)$ are a little
bit easier even without a boundary.
We are not going
to give complete proofs of all statements. Instead,
we concentrate on main ideas of the method. 
For the details an interested reader
can consult the original paper \cite{Gilkey1975}
and the monographs \cite{Gilkey:1995,Kirsten:2001wz}.

We start with very general properties of the heat
kernel coefficients. Let us consider 
a smooth compact Riemannian manifold $M$ without boundary.
To be able to define functions on $M$ which carry some discrete
(spin or gauge) indices  we need a vector bundle $V$ over $M$.
Let $D$ be an operator of Laplace type on $V$, 
and let $f$ be a smooth function on $M$.
There is an asymptotic expansion (\ref{asymptotex}) and

\noindent 1) Coefficients with odd index $k$ vanish,
$a_{2j+1}(f,D)=0$;

\noindent 2) Coefficients $a_{2j}(f,D)$ are locally
computable in terms of geometric invariants.

Already the existence of the power-law asymptotic expansion (\ref{asymptotex})
is a non-trivial statement. We postpone the discussion of this 
property to sec.\ \ref{s5other}.

The second statement above is very
important. It means that the heat kernel coefficients
can be expressed as integrals of local invariants:
\begin{equation}
a_k(f,D)= \ptr \iM \{ f(x) a_k (x;D) \} 
=\sum_I \ptr \iM \{ f u^I \mathcal{ A}_k^I(D) \} \,, \label{locinv}
\end{equation}
where $\mathcal{ A}_k^I$ are all possible independent invariants 
of the dimension $k$  constructed
from $E$, $\Omega$, $R_{\mu\nu\rho\sigma}$ and their derivatives.
We use usual assignments of the dimensions
when $E$ has dimension two, any derivative has dimension one, etc.
 $u^I$ are some constants. For example, if $k=2$ only two 
independent invariants exist. These are $E$ and $R$.
Note, that we can
always integrate by parts to remove all derivatives from $f$.
The first statement, $a_{2j+1}=0$, is clear now. 
One cannot construct an
odd-dimension invariant on a manifold without boundary.

Let us study further relations between the heat kernel coefficients
which will turn out to generate relations between the constants $u^I$.
Consider now the case when the manifold $M$ is a direct product
of two manifolds, $M_1$ and $M_2$, with coordinates $x_1$ and
$x_2$ respectively,
and the operator $D$ is a 
sum of two operators acting independently on $M_1$ and $M_2$,
$D=D_1\otimes 1 +1\otimes D_2$. This means that the bundle indices
are also independent. As an example, one can consider the vector Laplacian
on $M_1\times M_2$.  One can write symbolically
$\exp (-tD)=\exp (-tD_1)\otimes \exp (-tD_2)$. Next we multiply
both sides of this equation by $f(x_1,x_2)=f_1(x_1)f_2(x_2)$, take
the functional trace and perform the asymptotic expansion in $t$
to get
\begin{equation}
a_k(x;D)=\sum_{p+q=k} a_p(x_1;D_1)a_q(x_2;D_2)\,.
\label{product}
\end{equation}

The consequences of eq.\ (\ref{product})
are very far reaching. In particular, eq.\ (\ref{product}) allows to fix
the dependence of $u^I$ on the dimension of the manifold $M$.
Consider an even more specialised case when
one of the manifolds is a one-dimensional circle:  $M_1=S^1$, $0<x_1\le 2\pi$.
Let us make the simplest choice for $D_1$: 
$D_1=-\partial_{x_1}^2$. All geometric invariants associated with
$D$ are defined solely by the $D_2$-part. Moreover, all invariants are
independent of $x_1$.
Therefore, by the equation (\ref{locinv}),
\begin{eqnarray}
a_k(f(x_2),D)&=&\int_{S^1\times M_2}d^nx\sqrt{g}\sum_I
\ptr \{f(x_2) u^I_{(n)} 
\mathcal{ A}^I_k (D) \} \nonumber \\
&=&2\pi\int_{M_2}d^{n-1}x\sqrt{g}\sum_I
\ptr \{f(x_2) u^I_{(n)} \mathcal{ A}^I_k (D_2) \}\,.
\label{77hk}
\end{eqnarray}
Here dependence of the constants $u^I_{(n)}$ on the dimension $n$ of
the manifold is shown explicitly.
On the other hand, we can use (\ref{product}). Spectrum
of the operator $D_1$ is known. The eigenvalues are $l^2$,
$l \in {\mathbb {Z}}$. The heat kernel asymptotics for $D_1$
can be easily obtained
by using the Poisson summation formula
\begin{eqnarray}
K(t,D_1)&=&\sum_{l\in \mathbb{Z}} \exp (-tl^2)=\sqrt{\frac {\pi}t}
\sum_{l\in \mathbb{Z}} \exp (-\pi^2 l^2/t) \nonumber\\
&\simeq& \sqrt{\frac {\pi}t} +\mathcal{O}\left( e^{-1/t} \right)\,.
\label{eq:asPoi}
\end{eqnarray}
Since exponentially small terms have no effect on the heat kernel
coefficients, the only non-zero coefficient is
$a_0(1,D_1)=\sqrt{\pi}$. Therefore, we obtain by eq.\ (\ref{product})
\begin{equation}
a_k(f(x^2),D)={\sqrt \pi} \int_{M_2}d^{n-1}x\sqrt{g}\sum_I
\ptr
 \{f(x^2) u^I_{(n-1)} \mathcal{ A}^I_n (D_2) \} \,.  
\label{78hk}
\end{equation}
There are no restrictions on the 
operator $D_2$ or on the manifold $M_2$.
By comparing the equations (\ref{77hk}) and (\ref{78hk})
we obtain
\begin{equation}
u^I_{(n)}=\sqrt{4\pi} u^I_{(n+1)}. \label{uIn}
\end{equation}
This proves that the
constants $u^I$ depend on the dimension
$n$ only through the overall normalisation factor
$(4\pi )^{-n/2}$ .

To calculate the heat kernel coefficients we need also
the following variational equations (see \cite{Gilkey:1995},
Lemma 4.1.15):
\begin{eqnarray}
&&\frac{d}{d\epsilon} \vert_{\epsilon =0} a_k
(1, e^{-2\epsilon f} D)=(n-k) a_k (f,D)\,,\label{vara} \\
&&\frac{d}{d\epsilon} \vert_{\epsilon =0} a_k
(1,D-\epsilon F) =a_{k-2}(F,D)\,,\label{varb} \\
&&\frac{d}{d\epsilon} \vert_{\epsilon =0} a_{n-2}
(e^{-2\epsilon f}F,e^{-2\epsilon f} D)=0\,,\label{varc}
\end{eqnarray}
where $f$ and $F$ are some smooth functions.

To prove the first property (\ref{vara}) we note that
\begin{equation}
\frac{d}{d\epsilon} \vert_{\epsilon =0} {\rm Tr}
(\exp (-e^{-2\epsilon f}t D))=
{\rm Tr} (2ftD \exp (-tD)) =-2t\frac {d}{dt}
{\rm Tr} (f\exp (-tD))
\nonumber
\end{equation}
and expand both sides of this equation in the power series in $t$.
Eq.
(\ref{varb}) can be checked in a similar way. To prove (\ref{varc}) consider
the operator
\begin{equation}
D(\epsilon ,\delta )=e^{-2\epsilon f} (D -\delta F)\,.
\end{equation}
We use first (\ref{vara}) with $k=n$ to show:
\begin{equation}
0=\frac{d}{d\epsilon} \vert_{\epsilon =0} a_{n} (1,
D(\epsilon ,\delta )). \end{equation}
Then we vary the equation above with respect to $\delta$:
\begin{equation}
0= \frac{d}{d\delta}\vert_{\delta =0}
\frac{d}{d\epsilon} \vert_{\epsilon =0} a_{n} (1,
D(\epsilon ,\delta ))=\frac{d}{d\epsilon} \vert_{\epsilon =0}
\frac{d}{d\delta}\vert_{\delta =0} a_{n} (1,
D(\epsilon ,\delta ))\,. \end{equation}
Finally, eq.\ (\ref{varb}) yields (\ref{varc}).

Eq.\ (\ref{varb}) restricts dependence of the heat kernel coefficients
on the ``potential'' $E$ while (\ref{vara}) and (\ref{varc}) describe
properties of the heat kernel coefficients under local scale transformations.

To calculate the heat kernel coefficients we adopt the following
strategy. First we write down a general expression for $a_k$
containing all invariants $\mathcal{ A}_n^I$ of dimension $k$ with
arbitrary coefficients $u^I$. The constants $u^I$ are then
calculated by using the properties derived above. The first three
coefficients read \cite{Gilkey1975,Gilkey:1995}:
\begin{eqnarray}
&&a_{0}(f,D)=(4\pi)^{-n/2}\iM \ptr\{\alpha_0 f\} \label{a0gen} \\
&&a_2(f,D)
         =(4\pi)^{-n/2}\frac 16 \iM \ptr\{f(
\alpha_1E+\alpha_2R)\},\label{a2gen}\\
&&   a_4(f,D)=(4\pi)^{-n/2}\frac 1{360} \iM \ptr\{f(\alpha_3E_{;kk}
    +\alpha_4R E+\alpha_5E^2 \nonumber \\
&&\qquad\qquad
    +\alpha_6R_{;kk}+\alpha_7R^2+\alpha_8 R_{ij}R_{ij}
    +\alpha_9 R_{ijkl}R_{ijkl}+\alpha_{10}\Cur_{ ij}\Cur_{ij})\}.
\label{a4gen}
\end{eqnarray}
Instead of the $u^I$ we use rescaled constants $\alpha_I$. By (\ref{uIn})
the coefficients $\alpha_I$ are true constants, i.e. they do not depend on
$n$. One can check that indeed no more invariants exist. For example,
$R_{ij;ij}$ is proportional to $R_{;jj}$ due to the Bianchi identity.


The coefficient $\alpha_0$ follows immediately from the
heat kernel expansion for the ``free'' scalar Laplacian
on $S^1$ (see eq. (\ref{eq:asPoi})). We obtain
$\alpha_0=1$. Let us now use (\ref{varb}).
First take $k=2$. Then
\begin{equation}
\frac 16 \iM \ptr \left\{ \alpha_1 F \right\} =
\iM \ptr \{F\} \,. \label{4140}
\end{equation}
This gives $\alpha_1=6$. Take $k=4$ to see
\begin{equation}
\frac 1{360}\iM \ptr \{ \alpha_4 FR +2\alpha_5 FE \}
=\frac 16 \iM \ptr \{ \alpha_1 FE +\alpha_2 FR \}\,.
\label{4141}
\end{equation}
The equation (\ref{4141}) shows that $\alpha_5=180$,
$\alpha_4=60\alpha_2$. 

To proceed further we need local scale transformations defined in
(\ref{vara}) and (\ref{varc}). These scale transformations look
similar to the local Weyl transformations but are not exactly
the same. The scale transformations in (\ref{vara})
are designed in such a way that the operator $D$ always transforms
covariantly. This is not the case of the Weyl transformations of
an arbitrary operator of Laplace type. For example, the scalar Laplacian
(\ref{opDscalar}) is conformally covariant for a special value (\ref{confxi})
of the conformal coupling only.
Therefore, some of the basic quantities
are transformed in a somewhat unusual way. The metrics transforms
as $g_{\mu\nu}\to e^{2\epsilon f}g_{\mu\nu}$ thus defining standard
conformal properties of the Riemann tensor, the Ricci tensor and the
scalar curvature. The functions $a^\sigma$ and $b$ in eq. (\ref{D1})
transform homogeneously. Transformation properties of $\omega_\mu$
and $E$ are then defined through (\ref{oab}) and (\ref{Eab}).
One can obtain the following relations 
\cite{Gilkey1975,Gilkey:1995}
\begin{eqnarray}
&&\frac d{d\epsilon} \vert_{\epsilon =0} \sqrt{g}=nf\sqrt{g} \,,
\nonumber \\
&&\frac d{d\epsilon} \vert_{\epsilon =0} R_{ijkl}=-2fR_{ijkl}
+\delta_{jl}f_{;ik} +\delta_{;ik}f_{;jl}-\delta_{il}f_{;jk}
-\delta_{jk}f_{;il} \,,\nonumber \\
&&\frac d{d\epsilon} \vert_{\epsilon =0} E=-2fE +\frac 12 (n-2)f_{;ii}
\,,\nonumber \\
&&\frac d{d\epsilon} \vert_{\epsilon =0} R =-2fR -2(n-1)f_{;ii}
\,,\nonumber \\
&&\frac d{d\epsilon} \vert_{\epsilon =0} E_{;kk} =-4fE_{;kk}
-2f_{;kk}E +\frac 12 (n-2) f_{;iijj} +(n-6)f_{;k}E_{;k}
\,,\nonumber \\ 
&&\frac d{d\epsilon} \vert_{\epsilon =0} R E = -4f R E
+\frac 12 (n-2) f_{;ii}R -2(n-1)f_{;ii}E
\,,\nonumber \\
%\end{eqnarray}\begin{eqnarray}
&&\frac d{d\epsilon} \vert_{\epsilon =0} E^2 =-4fE^2 
+(n-2) f_{;ii}E
\,,\nonumber 
\end{eqnarray}\begin{eqnarray}
&&\frac d{d\epsilon} \vert_{\epsilon =0} R_{;kk} =-4fR_{;kk}
-2f_{;kk}R -2(n-1)f_{;iijj} +(n-6)f_{;i}R_{;i}
\,,\nonumber \\
&&\frac d{d\epsilon} \vert_{\epsilon =0} R^2 =-4fR^2
-4(n-1) f_{;ii} R
\,,\nonumber \\
&&\frac d{d\epsilon} \vert_{\epsilon =0} R_{ij}R_{ij} =-4fR_{ij}R_{ij}
-2f_{;ii}R -2(n-2)f_{;ij}R_{ij}
\,,\nonumber \\
&&\frac d{d\epsilon} \vert_{\epsilon =0} R_{ijkl}R_{ijkl} 
=-4fR_{ijkl}R_{ijkl} -8f_{;ij}R_{ij}
\,,\nonumber \\
&&\frac d{d\epsilon} \vert_{\epsilon =0} \Omega^2=-4f\Omega^2
\label{cvar1}
\end{eqnarray}
Let us remind that the indices $i,j,k,l$ are flat, so we can put them all
down and sum up over the repeated indices by contracting them with the
Kronecker $\delta$. 

Let us apply (\ref{varc}) to $n=4$.
\begin{equation}
\nonumber \frac{d}{d\epsilon} \vert_{\epsilon =0} a_2
(e^{-2\epsilon f}F,e^{-2\epsilon f} D)=0
\,. \end{equation}
By collecting the terms with $\ptr\iM \{ Ff_{;jj} \} $ we obtain
$\alpha_1=6\alpha_2$ and consequently $\alpha_2=1$ and
$\alpha_4=60\alpha_2=60$. This completes calculation of $a_2$.

Let $M=M_1\times M_2$ with a product metric and let $D=(-\Delta_1)+
(-\Delta_2)$ where $\Delta_{1,2}$ are scalar Laplacians on $M_1$ and
$M_2$ respectively\footnote{More precisely, we assume that in 
(\ref{opDscalar}) the potential $U=0$, the conformal coupling is minimal
$\xi=0$, and there is no gauge coupling $G_\mu=0$ on both $M_1$
and $M_2$.}.
Eq.\ (\ref{product}) yields
\begin{eqnarray}
&&a_4(1,-\Delta_1-\Delta_2)=
a_4(1,-\Delta_1)a_0(1,-\Delta_2)+
a_2(1,-\Delta_1)a_2(1,-\Delta_2)+ \nonumber \\
&&\qquad\qquad\qquad +a_0(1,-\Delta_1)a_4(1,-\Delta_2)
\label{prod4}
\end{eqnarray}
It is clear from (\ref{omEscal}) that $E=0$ and $\Omega =0$. 
By collecting the terms
with $R_1R_2$ (where $R_1$ and $R_2$ are scalar curvatures on
$M_1$ and $M_2$ respectively) we obtain
$$
\frac 2{360} \alpha_7 = \left( \frac{\alpha_2}6 \right)^2 \,.
$$
Consequently, $\alpha_7=5$. 

Let us apply (\ref{varc}) to $n=6$.
We obtain with the help of the variational equations (\ref{cvar1})
\begin{eqnarray}
&&0=\ptr \iM \{ F( (-2\alpha_3-10\alpha_4+4\alpha_5)f_{;kk}E
\nonumber\\
&&\qquad +(2\alpha_3 -10\alpha_6)f_{;iijj}
\nonumber\\
&&\qquad +(2\alpha_4-2\alpha_6-20\alpha_7-2\alpha_8 )f_{;ii}R
\nonumber \\
&&\qquad +(-8\alpha_8-8\alpha_9 )f_{;ij}R_{ij} )\} \,.
\label{var1012}
\end{eqnarray}
The coefficients in front of independent invariants in (\ref{var1012})
must be zero. We determine $\alpha_3=60$, $\alpha_6=12$,
$\alpha_8=-2$ and $\alpha_9=2$.

The most elegant way to calculate the remaining constant $\alpha_{10}$
is based on the Gauss-Bonnet theorem 
(see \cite{Gilkey:1995}). We use here a more lengthy way 
\cite{Nepomechie:1985wt} which
however works perfectly on flat manifolds without boundary and
of trivial topology of $\mathbb{R}^n$. 
Note that $\mathbb{R}^n$ is non-compact.
To make
the heat kernel well defined we should suppose certain fall-off
conditions on the background fields and on the smearing function $f$.
The basis in the space of the square integrable
functions is given by the plane waves $\exp (ikx)$.
Therefore, for $M=\mathbb{R}^n$ with flat metric the heat kernel
reads:
\begin{eqnarray}
&&K(f;t)={\rm Tr}_{L^2} (f\exp (-tD))\nonumber \\
&&\qquad =\int d^nx \int 
\frac {d^nk}{(2\pi)^n} e^{-ikx}\ptr \{f(x) \exp (-tD) e^{ikx}\}
\nonumber \\
&&\qquad =\int d^nx \int \frac {d^nk}{(2\pi)^n} \ptr\{ f(x) 
\exp (t((\nabla^\mu +ik^\mu )^2+E))\} \,.
\label{flathk1}
\end{eqnarray}
The following integrals will be useful:
\begin{eqnarray}
&&\int \frac {d^nk}{(2\pi)^n} e^{-tk^2} =\frac 1{(4t\pi)^{n/2}}
\, , \nonumber \\
&&\int \frac {d^nk}{(2\pi)^{n}}  e^{-tk^2} k^\mu k^\nu
=\frac 1{(4t\pi)^{n/2}}\frac 1{2t} g^{\mu \nu} , \nonumber \\
&&\int \frac {d^nk}{(2\pi)^n} e^{-tk^2}  k^\mu k^\nu k^\rho k^\sigma
=\frac 1{(4t\pi)^{n/2}} \frac 1{4t^2} (g^{\mu\nu}g^{\rho \sigma}+
g^{\mu\rho}g^{\nu\sigma}+g^{\mu\sigma}g^{\nu\rho} ) .\label{eq:ex2}
\end{eqnarray}
Now we isolate $\exp (-tk^2)$ on the right hand side of (\ref{flathk1})
and expand the rest of the exponent in a power series of $t$.
\begin{eqnarray}
K(f;t) &=&\int d^nx \int \frac {d^4k}{(2\pi)^n} e^{-tk^2}
\ptr \left\{ f(x)\left( 1+t(\nabla^2 +E)
-\frac {t^2}2 4(k\nabla ) \right. \right. \nonumber \\
&+& \frac {t^2}2 (\nabla^2 \nabla^2 +\nabla^2 E +E\nabla^2 +E^2)
\nonumber \\
&-& \frac {4t^3}6 ((k\nabla )^2E+E(k\nabla )^2 +(k\nabla )E(k\nabla ))
\nonumber \\
&-& \frac {4t^3}6
((k\nabla )^2\nabla^2+\nabla^2(k\nabla )^2 +(k\nabla )\nabla^2(k\nabla ))
\nonumber \\
&+&\left. \left. \frac {16t^4}{24} (k\nabla )^4 +\dots \right) \right\}
\end{eqnarray}
We use the integrals (\ref{eq:ex2}) to obtain
\begin{eqnarray}
K(f;t)&=& \frac 1{(4\pi t )^{n/2}}   \int d^nx 
  \ptr \left\{ f(x) \left( 1 + t E \right.\right.
\nonumber\\
&+&\frac {t^2}2 (\nabla^2 \nabla^2 +\nabla^2 E+E\nabla^2+E^2)
\nonumber \\
&-& \frac {t^2}3 (\nabla^2 E +E\nabla^2 +\nabla^\mu E\nabla_\mu )
\nonumber \\
&-& \frac {t^2}3 (2 \nabla^2 \nabla^2 +\nabla^\mu \nabla^2 \nabla_\mu )
\nonumber \\
&+& \left. \left.
\frac {t^2}6 (\nabla^\mu \nabla^\nu \nabla_\mu \nabla_\nu
+\nabla^2 \nabla^2+\nabla^\mu \nabla^2 \nabla_\mu )+O(t^3)
\right) \right\} \,.
\end{eqnarray}
All derivatives combine into commutators. Finally we get:

\begin{eqnarray}
&&K(f;t)=\frac 1{(4\pi t)^{n/2}} \int d^nx \ptr \left\{ f(x)
\left( 1 +t E \right. \right. \nonumber \\
&&\qquad\qquad\qquad \left. \left. + t^2 \left(
\frac 12 E^2 +\frac 16 E_{;\mu\mu}+
\frac 1{12} \Omega_{\mu \nu}\Omega^{\mu \nu} \right) +O(t^3)\right)
\right\}\,.
\label{eq:ex3}
\end{eqnarray}
From the equation (\ref{eq:ex3}) we conclude that $\alpha_{10}=30$.
We have obtained also an independent confirmation for the values 
of $\alpha_0$, $\alpha_1$, $\alpha_3$ and $\alpha_5$.

The method we used above in eqs. (\ref{flathk1}) - (\ref{eq:ex3})
can be applied to more general manifolds and operators as well.
The key ingredient is a convenient basis which should be used
instead of the plane waves.  
For the case of a box with (anti-) periodic boundary condition
such a basis is rather obvious (see  
\cite{Andrianov:1984fg,Andrianov:1984qj}). On curved manifold
$M$ one has to use the so-called geodesic waves 
\cite{Fujikawa:1979ay,Fujikawa:1980eg,Yajima:1988pj,Yajima:1996jk,
Novozhilov:1991nm}
although calculations with ordinary plane 
waves are also possible \cite{Ceresole:1989hn}.

For calculation of the coefficient $a_6$ we refer to
\cite{Gilkey1975}. The results for the leading heat kernel
coefficients are summarised in the following equations
\begin{eqnarray}
&&a_{0}(f,D)=(4\pi)^{-n/2}\iM \ptr\{f\}. \label{a0nobou}\\
&&a_2(f,D)
         =(4\pi)^{-n/2}6^{-1}\iM \ptr\{f(6E+R)\}.\label{a2nobou}\\
&&   a_4(f,D)=(4\pi)^{-n/2}360^{-1}\iM \ptr\{f(60E_{;kk}
    +60R E+180E^2 \nonumber \\
&&\qquad\qquad
    +12R_{;kk}+5R^2-2 R_{ij}R_{ij}
    +2 R_{ijkl}R_{ijkl}+30\Cur_{ ij}\Cur_{ij})\}.\label{a4nobou} \\
&&a_{6}(f,D)=(4\pi)^{-n/2}\iM \ptr\bigl\{ \frac f{7!}(
    18R_{;ii jj}+17R_{;k}R_{;k}
    -2R_{ij;k}R_{ ij;k} \nonumber \\
&&\qquad\qquad -4R_{jk;n}R_{jn;k}
    +9R_{ij kl;n}R_{ij kl;n}+28RR_{;nn}
    -8R_{jk}R_{jk;nn} \nonumber \\
&&\qquad\qquad
    +24R_{ jk}R_{jn;kn}
    +12R_{ij kl}R_{ij kl ;nn}
    +35/9R^{3}
    -14/3RR_{ij}R_{ij} \nonumber \\
&&\qquad\qquad
    +14/3R R_{ijkl}R_{ijkl}
     -208/9R_{jk}R_{jn}R_{kn}
     -64/3R_{ ij}R_{kl}R_{ik jl} \nonumber \\
&&\qquad\qquad
     -16/3R_{jk}R_{jn l i}R_{kn l i}
     -44/9R_{ij kn}R_{ij l p}R_{kn l p} \nonumber\\
&&\qquad\qquad
     -80/9R_{ij kn}R_{il  kp}R_{jl  np})
     +360^{-1}f(
      8\Cur_{ij;k} \Cur_{ij;k}
     +2\Cur_{ij;j}\Cur_{ ik;k}\nonumber \\
&&\qquad\qquad
     +12\Cur_{ij;kk}\Cur_{ij}
     -12\Cur_{ij}\Cur_{jk}\Cur_{ki}
     -6R_{ij kn}\Cur_{ij}\Cur_{kn}
     -4R_{jk}\Cur_{jn}\Cur_{kn}
\nonumber \\
&&\qquad\qquad
       +5R\Cur_{kn}\Cur_{kn}
     +6E_{;ii jj}+60EE_{;ii}
      +30E_{;i}E_{;i}
     +60E^{3}
\nonumber \\
&&\qquad\qquad
     +30E\Cur_{ij}\Cur_{ij}
     +10R E_{;kk}+4R_{jk}E_{;jk}
     +12R_{;k}E_{;k}+30EER
     \nonumber \\
&&\qquad\qquad
      +12ER_{;kk}+5ER^2
     -2ER_{ij}R_{ij}+2ER_{ijkl}R_{ijkl})\}.\label{a6nobou}
\end{eqnarray}

Everyone who ever attempted calculations of the heat kernel
coefficients on curved background for arbitrary spin should
appreciate that the method presented here is a quite
efficient one. One should also take into account that some
of the universal constant were indeed calculated twice.

The coefficients $a_0$ -- $a_4$ are contained in 
\cite{DeWitt:1965,McKean:1967}.
$a_6$ was first computed by Gilkey \cite{Gilkey1975}.
The next coefficient $a_8$ has been calculated by Amsterdamski
{\it et al} \cite{Amsterdamski:1989bt} for the scalar Laplacian
and by Avramidi \cite{Avramidi:1990ug,Avramidi:1991je} for
the general operator of Laplace type.
The coefficient $a_{10}$ has been calculated 
by van de Ven\cite{vandeVen:1997pf}.
Higher heat kernel coefficients in flat space were studied in
\cite{Fliegner:1995zc,Fliegner:1998rk}.
%%%%%%%%%%%%%%%%
\subsection{Examples}\label{s4examples}
Here we consider several simple physical systems in four dimensions
and calculate the heat kernel coefficient $a_4$ which defines the
one-loop divergences in the zeta function regularization. 
%%%%
\subsubsection{Yang--Mills theory in flat space}\label{s4YM}
Our first example is pure Yang--Mills theory in flat space. We are 
interested in the ``total'' heat kernel coefficient $a_4^{\rm tot}$
defined by (\ref{atot}). Let us start with the first term describing
contribution from the vector fields. We choose the gauge (\ref{bg-gauge}).
The only non-vanishing invariants are $E$ and $\Omega$. The coefficient
$a_4$ is quadratic in these quantities. By using eqs.\ (\ref{EYM})
and (\ref{OmegaYM}) we obtain:
\begin{eqnarray}
&&\ptr \left( E^2\right)=E_{\nu\beta}^{\alpha\rho}E_{\rho\alpha}^{\beta\nu}=
4 F_{\rho\nu}^\delta F_{\rho\nu}^\gamma K_{\delta\gamma} \,,
\label{trE2vect} \\
&&\ptr \left( \Omega_{ij} \Omega_{ij} \right) = 
-4 F_{\rho\nu}^\delta F_{\rho\nu}^\gamma K_{\delta\gamma} \,,
\label{trO2vect}
\end{eqnarray} 
where 
\begin{equation}
K_{\delta\gamma}=c_{\alpha\beta}^{\delta}c_{\alpha\beta}^{\gamma} 
\label{Killform}
\end{equation}
is the Killing form of the gauge algebra.

Let us remind that the indices $i,j,k$ in (\ref{a4nobou}) refer to a local
orthonormal frame. In flat space they may be identified with the
vector indices $\mu,\nu,\rho$. In the case of vector fields the trace
in (\ref{trE2vect}) and (\ref{trO2vect}) is taken over {\it pairs}
consisting of a gauge index $\alpha$, $\beta$ or $\gamma$ and of a
vector index $\mu$, $\nu$ or $\rho$.

For the ghost operator we have to put $R=U=0$ in (\ref{omEscal})
and identify the connection with the background gauge field
according to (\ref{ghostcon}). Consequently, $E^{[{\rm gh}]}=0$ and
\begin{equation}
\ptr \left( \Omega_{ij}^{[{\rm gh}]} \Omega_{ij}^{[{\rm gh}]} \right)=
-F_{\rho\nu}^\delta F_{\rho\nu}^\gamma K_{\delta\gamma} \,.\label{trO2ghost}
\end{equation}
Next we substitute (\ref{trE2vect}), (\ref{trO2vect}) and 
(\ref{trO2ghost}) in (\ref{a4nobou}) to obtain
\begin{equation}
a_4^{[{\rm tot}]}=a_4^{[{\rm vec}]}-2a_4^{[{\rm gh}]}=
\frac{11}{96\pi^2} \int_M d^4x\, \sqrt{g}\, 
F_{\rho\nu}^\delta F_{\rho\nu}^\gamma K_{\delta\gamma}\,.
\label{YMa4}
\end{equation}

The Yang--Mills gauge group $\mathcal{G}$ has usually a direct product
structure, $\mathcal{G}=\mathcal{G}_1\times \mathcal{G}_2 \times\dots$,
so that on each of the irreducible components $\mathcal{G}_i$ the Killing
form $K$ is proportional to the unit matrix. Therefore, the one-loop
divergence (\ref{YMa4}) reproduces the structure of the classical action
with different charges for each $\mathcal{G}_i$. We also recover the
coefficient $11/3$ which is familiar from computations of
the Yang--Mills beta functions.
%%%
\subsubsection{Free fields in curved space}\label{s4free}
Consider free quantum fields on a Riemannian manifold without boundaries.
``Free'' means that we neglect all interactions except for the one
with the background geometry. The heat kernel coefficients can be
expressed then in terms of local invariants of the metric.
In particular,
\begin{equation}
a_4(x)=\frac 1{2880\pi^2} \left[
aC_{\mu\nu\rho\sigma}C^{\mu\nu\rho\sigma} +b\left( R_{\mu\nu}R^{\mu\nu}
-\frac 13 R^2 \right) +c {R_{;\mu}}^\mu +d R^2 \right]\,,
\label{a4in4dim}
\end{equation} 
where $C_{\mu\nu\rho\sigma}$ is the Weyl tensor,
\begin{equation}
C_{\mu\nu\rho\sigma}C^{\mu\nu\rho\sigma}=R_{\mu\nu\rho\sigma}
R^{\mu\nu\rho\sigma} - 2R_{\mu\nu}R^{\mu\nu} +\frac 13 R^2 \,.
\label{Weylsquare}
\end{equation}
$a$, $b$, $c$ and $d$ are some constants depending on the spin.
The first two structures in (\ref{a4in4dim}) which appear together
with $a$ and $b$ are conformally invariant in four dimensions.
This explains our choice of the basis in the space of invariants.

The constants $a$, $b$, $c$ and $d$ can be evaluated by substituting
particular expressions for $E$ and $\Omega$ obtained in sec.\ 
\ref{sec3} in (\ref{a4nobou}). Alternatively, one can use the
analysis of Christensen and Duff \cite{Christensen:1979md}
who calculated $a_4$ for arbitrary spin fields (see also
\cite{Birrell:1982ix}). The results are collected in Table \ref{a4table}.
Some comments are in order. Spin $1/2$ means 4-component Dirac spinors.
For spin 1 and spin 2 fields we took into
account contributions from corresponding ghost fields. Note, that
vector Yang-Mills fields and vector ghosts for gravity interact
{\em differently} with the background geometry. The cosmological
constant is taken to be zero. 
\begin{table}%\label{a4table}
\caption{$a_4$ for various spins}
\begin{tabular}{|l|c|c|c|c|}\hline \label{a4table}
Spin & $a$ & $b$ & $c$ & $d$ \\ \hline
$0$ &    $1$ & $1$ & $30\xi -6$ & $90 (\xi -1/6)^2$ \\
$1/2$&    $-7/2$ & $-11$ & $6$ & $0$ \\
$1$ &    $-13$  & $62$  & $18$ & $0$ \\
$2$ &    $212$  & $0$   & $0$  & $717/4$ \\ \hline %\label{a4table}
\end{tabular}
\end{table}  

As a physical application of the heat kernel expansion in curved
space we may mention, for example, the asymptotic conformal invariance
phenomenon which was studied by using this technique in
\cite{Buchbinder:1984zx,Buchbinder:1985ba,Buchbinder:1984gj,Buchbinder:1988pm}.

A similar analysis can be performed also in the presence of boundaries
\cite{Moss:1994jj}.
%%%%%%%%%%%%%%%%
\subsection{DeWitt iterative procedure}\label{s4iter}
The iterative method by DeWitt uses separation of the heat kernel
with non-coinciding arguments into a leading part (which is non-analytic
in $t$) and the power-law corrections. For a flat manifold this separation
has been described in sec.\ \ref{sec1} (cf.\ eqs.\ (\ref{simplestHK})
and (\ref{sHKexp})). Let us consider a massless scalar field on
a curved compact Riemannian manifold $M$. The DeWitt ansatz reads
in this case:
\begin{equation}
K(t;x,y;D)=(4\pi t)^{-n/2} \Delta_{VVM}^{1/2}(x,y) \exp 
\left( -\frac{\sigma(x,y)}{2t} \right) \Xi (t;x,y;D)\,,
\label{DWansatz}
\end{equation}
where $\sigma (x,y)$ is one half the square of the length of
the geodesic connecting $x$ and $y$. In Cartesian coordinates
on a flat manifold $\sigma_{\rm flat}(x,y)=\frac 12 (x-y)^2$.
$\Delta_{VVM}$ is the so called Van Vleck--Morette determinant
\begin{equation}
\Delta_{VVM}(x,y) =\frac{\det \left( -\frac{\partial}{\partial x^\mu}
\frac{\partial}{\partial y^\nu} \sigma (x,y) \right)}{\sqrt{g(x)g(y)}}
\,.\label{VVMdet}
\end{equation}

As a consequence of the heat equation (\ref{heateq}) the kernel
$\Xi$ should satisfy
\begin{equation}
\left( \partial_t +t^{-1} (\nabla^\mu \sigma )\nabla_\mu +
\Delta_{VVM}^{-1/2} D \Delta_{VVM}^{1/2} \right) \Xi =0
\label{eqforXi}
\end{equation}
with the initial condition
\begin{equation}
\Xi (0;x,y;D)=1 \,.\label{initXi}
\end{equation}

The essence of the DeWitt method is to look for a solution of the
equation (\ref{eqforXi}) in the following form
\begin{equation}
\Xi (t;x,y;D) = \sum_{j=0}^\infty t^{j}b_{2j}(x,y;D)\,.
\label{ansatzforXi}
\end{equation}
The initial condition (\ref{initXi}) yields
\begin{equation}
b_0=1\,.\label{b0DW}
\end{equation}
The recursion relation 
\begin{equation}
\left( j +(\nabla^\mu \sigma )\nabla_\mu \right) b_{2j}+
\Delta_{VVM}^{-1/2} D \Delta_{VVM}^{1/2} b_{2(j-1)}=0\,,
\label{recursionDW}
\end{equation}
which follows from (\ref{eqforXi}) and (\ref{ansatzforXi}),
allows, in principle, to find the coincidence limits $x=y$
of higher heat kernel coefficients $b_{2j}$. An important
ingredient of such calculations is the coincidence limits
of symmetrized derivatives of the geodesic interval $\sigma (x,y)$
\cite{Synge:1931,Barvinsky:1985an}. This method becomes
very cumbersome beyond $a_4$. A refined nonrecursive procedure
to solve the DeWitt equation (\ref{recursionDW}) was used
by Avramidi \cite{Avramidi:1991je} to calculate $a_8$
(see also \cite{Avramidi:2001ns} for a short overview).
The method of DeWitt can be naturally extended to treat
coincidence limits of the derivatives of $b_{2j}(x,y;D)$.

The recursion relations can be generalised to the case
of manifolds with boundaries \cite{McAvity:1991we,McAvity:1992rf}.
However, for a practical use the functorial methods of sec.\ 
{\ref{sec5}} seem to be more convenient.
%%%%%%%%%%%%%%%
\subsection{Non-minimal operators}\label{s4nonmin}
Quantisation of gauge theories oftenly leads to second order
differential operators which are not of the Laplace type.
For example, by taking $\kappa\ne 1$ in (\ref{Svgf}) 
one obtains the following operator acting on the gauge field
fluctuations:
\begin{equation}
D^{[{\rm nm}]}_{\mu\nu} =-\Delta g_{\mu\nu} +
\left( 1-\frac 1\kappa \right) \nabla_\nu\nabla_\mu \,,
\label{nonminop}
\end{equation}
where, for simplicity, we suppose that the manifold $M$ is flat
($R_{\mu\nu\rho\sigma}=0$) and the gauge group is abelian
($c_{\alpha\beta}^\gamma=0$). The leading symbol of the operator
(\ref{nonminop}) (the part with the highest derivatives) has a 
non-trivial matrix structure. Such operators are called {\it 
non-minimal}.

In some simple cases (see, e.g., 
\cite{GBF:1991,Alexandrov:1996gu,Vassilevich:1995bg}) the spectral
problem for non-minimal operators can be reduced to the Laplacians.
If the leading part of a non-minimal operator has a form similar to
(\ref{nonminop}) but the lower order part is more or less 
arbitrary\footnote{This case covers most of the physical applications.},
necessary generalisations of the DeWitt technique were suggested by
Barvinsky and Vilkovisky \cite{Barvinsky:1985an}. The technique was
further developed in \cite{Gusynin:1991ek,Cho:1995ic}. Complete
calculation of $a_4$ required extensive use of the computer algebra
\cite{Gusynin:1997qs}. Most general non-minimal operators were
considered in \cite{Avramidi:2001tx} where, because of great technical
complexity, only the first two heat kernel coefficients were analysed.
Some calculations in various physical systems with non-minimal operators
can be found in 
\cite{Fradkin:1982ts,Parker:1984pe,Guendelman:1994ke,Pronin:1997bu}.
